% ---------------------------------------------------------------------
% 2.8 Redes neuronales artificiales y aprendizaje profundo
% ---------------------------------------------------------------------
\section[REDES NEURONALES ARTIFICIALES Y APRENDIZAJE PROFUNDO]{Redes neuronales artificiales y
aprendizaje profundo}

\subsection{La neurona artificial}

Las redes neuronales artificiales son modelos inspirados en el sistema nervioso, en los que las
unidades de cómputo, llamadas neuronas, se conectan entre sí mediante pesos que cumplen un papel
análogo al de la fuerza de las conexiones sinápticas en los organismos biológicos
\cite{aggarwal2018}. Cada neurona recibe un vector de entradas $\mathbf{x} = (x_1, \dots, x_n)$,
las multiplica por sus pesos $w_i$, suma el resultado junto con un sesgo $b$ y aplica una función
de activación $\varphi$ no lineal \cite{delbrio2002}:
\begin{equation}
y = \varphi\!\left(\sum_{i=1}^{n} w_i x_i + b\right) = \varphi\!\left(\mathbf{w}^{\top}\mathbf{x} + b\right).
\label{eq:neurona}
\end{equation}
Al organizar muchas de estas unidades en capas, de forma que
la salida de una capa sea la entrada de la siguiente, se obtiene el perceptrón multicapa. Las capas
intermedias se llaman ocultas, y el número de capas se conoce como profundidad de la red. De ahí
proviene el término aprendizaje profundo \cite{goodfellow2016}.

\subsection{Funciones de activación}

Sin la función de activación, una red de cualquier profundidad equivaldría a una sola
transformación lineal, y es gracias a ella que la red puede representar relaciones no
lineales \cite{goodfellow2016}. La Tabla~\ref{tab:activaciones} presenta las funciones más
utilizadas y su uso en las arquitecturas para la detección de la P300.

\begin{table}[htbp]
\centering
\small
\renewcommand{\arraystretch}{1.8}
\begin{tabularx}{\textwidth}{|L{2.6cm}|L{4.3cm}|L{1.6cm}|Y|}
\hline
\textbf{Función} & \textbf{Definición} & \textbf{Rango} & \textbf{Uso típico} \\ \hline
Sigmoide & $\sigma(z) = \dfrac{1}{1+e^{-z}}$ & $(0,1)$ & Salida de clasificación binaria (probabilidad de P300). \\[8pt] \hline
Tangente hiperbólica & $\tanh(z) = \dfrac{e^{z}-e^{-z}}{e^{z}+e^{-z}}$ & $(-1,1)$ & Capas ocultas en \cite{cecotti2011,alvarado2021}. \\[8pt] \hline
ReLU & $\max(0,z)$ & $[0,\infty)$ & Capas ocultas de redes profundas. \\[4pt] \hline
ELU & $z$ si $z>0$, y $\alpha(e^{z}-1)$ si $z\le0$ & $(-\alpha,\infty)$ & Capas ocultas de EEGNet \cite{lawhern2018}. \\[4pt] \hline
Softmax & $\dfrac{e^{z_k}}{\sum_{j} e^{z_j}}$ & $(0,1)$, suma 1 & Salida de clasificación multiclase. \\ \hline
\end{tabularx}
\caption{Funciones de activación más utilizadas. Elaboración propia con información de \cite{goodfellow2016,aggarwal2018}.}
\label{tab:activaciones}
\end{table}

\subsection{Función de pérdida}

El entrenamiento de una red consiste en ajustar sus pesos para minimizar una función de pérdida
que mide la diferencia entre la salida de la red y la etiqueta esperada. Para la clasificación
binaria, con etiquetas $y_i\in\{0,1\}$ y salidas $\hat{y}_i\in(0,1)$, se utiliza la entropía
cruzada binaria \cite{goodfellow2016}:
\begin{equation}
\mathcal{L}_{\text{BCE}} = -\frac{1}{M}\sum_{i=1}^{M}\Big[\,y_i\log\hat{y}_i + (1-y_i)\log(1-\hat{y}_i)\Big].
\label{eq:bce}
\end{equation}
Debido al desbalance de las épocas del P300 Speller, en el que por cada época positiva hay cinco
negativas, es común ponderar el primer término con un factor $\alpha>1$ para que los errores sobre
la clase minoritaria tengan mayor peso.

\subsection{Descenso de gradiente y retropropagación}

Para minimizar la pérdida se utiliza el método de descenso de gradiente, que actualiza los pesos
en la dirección opuesta al gradiente de la función de pérdida,
\begin{equation}
\mathbf{w}^{(t+1)} = \mathbf{w}^{(t)} - \eta\, \nabla_{\mathbf{w}}\mathcal{L}\big(\mathbf{w}^{(t)}\big),
\label{eq:gradiente}
\end{equation}
donde $\eta$ es la tasa de aprendizaje. Si es muy pequeña, el entrenamiento resulta lento, y si es
muy grande, los pesos oscilan alrededor del mínimo sin alcanzarlo \cite{goodfellow2016}. En redes de
varias capas, este gradiente se calcula con el algoritmo de retropropagación, que propaga el error
desde la capa de salida hacia las capas anteriores \cite{rumelhart1986}.

\subsection{Regularización}

Dado que las redes profundas tienen una gran cantidad de parámetros y los conjuntos de datos de EEG
son relativamente pequeños, es indispensable emplear técnicas de regularización para evitar el
sobreajuste. Entre las más utilizadas se encuentran el \textit{dropout}, que desactiva
aleatoriamente una fracción de las neuronas en cada iteración de entrenamiento
\cite{srivastava2014}, la normalización por lotes, que normaliza las activaciones de cada capa y
estabiliza el entrenamiento \cite{ioffe2015}, y la parada temprana, que detiene el entrenamiento
cuando la pérdida de validación deja de mejorar \cite{goodfellow2016}. Las dos primeras forman
parte de arquitecturas para la detección de la P300 como EEGNet \cite{lawhern2018}.

